\documentclass[12pt,a4paper]{article}

% --- Mathématiques de précision ---
\usepackage{amsmath, amssymb, esvect, mathtools}
\newcommand{\R}{\mathbb{R}}
\newcommand{\N}{\mathbb{N}}
\newcommand{\e}{\mathrm{e}}
\newcommand{\vect}[1]{\vv{#1}}

% --- Encodage et typographie française ---
\usepackage[utf8]{inputenc}
\usepackage{fontspec}
\usepackage{libertinus-otf}
\usepackage[french]{babel}
\usepackage{microtype}
\usepackage{icomma}

% --- Suppression de l'indentation ---
\setlength{\parindent}{0pt}

% --- Géométrie calibrée (exactement 1 page par exercice) ---
\usepackage[
    top=1.6cm,
    bottom=1.6cm,
    left=1.6cm,
    right=1.6cm,
    headheight=15pt,
    headsep=12pt,
    footskip=22pt
]{geometry}

\usepackage{fancyhdr}
\usepackage{lastpage}
\usepackage{tabularx}
\usepackage{tcolorbox}
\usepackage{enumitem}

% --- Filigrane et graphisme ---
\usepackage{eso-pic}
\usepackage{tikz}
\usepackage{xcolor}

% --- Réglage des espacements pour affichages mathématiques ---
\setlength{\abovedisplayskip}{6pt}
\setlength{\belowdisplayskip}{6pt}
\setlength{\abovedisplayshortskip}{3pt}
\setlength{\belowdisplayshortskip}{3pt}

% --- Colonne centrée pour tableau homogène ---
\newcolumntype{Y}{>{\centering\arraybackslash}X}

% --- Commandes d'exercices aérées ---
\newcommand{\exercice}[3][5]{%
    \noindent{\Large\textbf{Exercice #2}} \hfill \textbf{#1 points}\\[0.20cm]%
    {\normalsize\textit{#3}}%
    \par\vspace{0.28cm}\hrule\vspace{0.55cm}%
}
\newcommand{\partie}[1]{\vspace{0.45cm}\noindent\textbf{\large Partie #1}\par\vspace{0.25cm}}

% --- Listes numérotées : 1. et a. TOUS EN GRAS ---
\setlist[enumerate,1]{label=\textbf{\arabic*.}, leftmargin=*, itemsep=0.35em, parsep=0pt, topsep=0.2em}
\setlist[enumerate,2]{label=\textbf{\alph*.}, leftmargin=*, itemsep=0.25em, parsep=0pt, topsep=0.15em}

% --- SQUELETTE : marqueur magenta des zones à remplir (utilisable en texte ET en mode maths) ---
\newcommand{\ph}[1]{\ifmmode\text{\color{magenta}[#1]}\else\textcolor{magenta}{[#1]}\fi}

% --- En-têtes et pieds de page pour le Sujet 1 ---
\pagestyle{fancy}
\fancyhf{}
\fancyhead[L]{\small\textit{Bac Blanc -- Sujet 1}}
\fancyhead[C]{\small\textit{Sujet 1 (Jour 1)}}
\fancyhead[R]{\small\textbf{Term\textsuperscript{spé}}}
\fancyfoot[C]{\thepage\ / \pageref{LastPage}}
\renewcommand{\headrulewidth}{0.4pt}
\renewcommand{\footrulewidth}{0.4pt}

% =========================================================================
\begin{document}

% -------------------------------------------------------------------------
% PAGE DE GARDE OFFICIELLE (Page 1)
% -------------------------------------------------------------------------
\thispagestyle{empty}

% FILIGRANE DISCRET ET ÉLÉGANT
\AddToShipoutPictureFG*{%
    \begin{tikzpicture}[remember picture, overlay]
        \node[
            rotate=32,
            scale=1.1,
            text=black,
            opacity=0.06,
            font=\fontsize{72pt}{72pt}\bfseries\sffamily
        ] at (current page.center) {SP{\'E}CIALIT{\'E}};
    \end{tikzpicture}%
}

% En-tête de la page de garde
\noindent
\begin{minipage}[t]{0.52\textwidth}
    \textbf{\Large LYCÉE JEAN ROSTAND}\\[0.12cm]
    \textsc{Année scolaire 2026--2027}\\
    \textit{Classe de Terminale Générale}
\end{minipage}
\hfill
\begin{minipage}[t]{0.45\textwidth}
    \raggedleft
    \textbf{\large BACCALAURÉAT BLANC}\\[0.10cm]
    \textbf{Épreuve de spécialité}\\[0.10cm]
    \textit{Sujet 1 --- Jour 1}\\[0.12cm]
    \small Date : \ph{JJ mois AAAA}
\end{minipage}

\vspace{0.30cm}
\hrule height 0.8pt
\vspace{0.45cm}

\vfill

% Titre officiel
\begin{center}
    {\LARGE \textbf{BACCALAURÉAT BLANC}}\\[0.20cm]
    {\large \textbf{SESSION \ph{MOIS AAAA}}}\\[0.15cm]
    {\textbf{\Large SUJET 1 --- JOUR 1}}\\[0.45cm]

    \begin{tcolorbox}[colback=white, colframe=black, arc=0mm, boxrule=1.5pt, width=0.96\textwidth, top=4mm, bottom=4mm]
        \centering
        {\fontsize{26pt}{30pt}\selectfont\textbf{MATHÉMATIQUES}}\\[0.30cm]
        {\Large \textsc{Épreuve d'Enseignement de Spécialité}}
    \end{tcolorbox}

    \vspace{0.50cm}
    {\large \textbf{Durée de l'épreuve : 4 heures}}\\[0.25cm]
    \textbf{L'usage de la calculatrice avec \underline{mode examen actif} est autorisé.}\\[0.10cm]
    \textit{L'usage de la calculatrice sans mémoire, « type collège », est autorisé.}
\end{center}

\vfill

% Tableau de structure officiel : 4 / 5 / 6 / 5 = 20 points
\begin{center}
\renewcommand{\arraystretch}{1.5}
\begin{tabularx}{0.96\textwidth}{|>{\centering\arraybackslash\bfseries}m{2.6cm}|*{4}{Y|}>{\centering\arraybackslash\bfseries}m{2.2cm}|}
    \hline
    Exercice & \textbf{Exercice 1} & \textbf{Exercice 2} & \textbf{Exercice 3} & \textbf{Exercice 4} & \textbf{Total} \\ \hline
    Thème    & \ph{Thème 1} & \ph{Thème 2} & \ph{Thème 3} & \ph{Thème 4} & --- \\ \hline
    Barème   & 4 points & 5 points & 6 points & 5 points & 20 points \\ \hline
\end{tabularx}
\end{center}

\vfill

% MESSAGE SURVEILLANTS
\begin{center}
\begin{tcolorbox}[colback=gray!5, colframe=black, boxrule=1pt, arc=1mm, width=0.96\textwidth, top=3mm, bottom=3mm, left=4mm, right=4mm]
    \centering
    \textbf{\textsc{Consigne importante aux surveillants}}\\[0.15cm]
    \normalsize
    Nous remercions les surveillants de veiller à ce que chaque candidat inscrive lisiblement sur sa copie le \textbf{nom de son professeur de spécialité}.
\end{tcolorbox}
\end{center}

\vfill

% Consignes officielles
\begin{tcolorbox}[colback=white, colframe=black, arc=0mm, boxrule=0.8pt, title=\textbf{\textsc{Consignes aux candidats}}, coltitle=black, colbacktitle=gray!15, top=3mm, bottom=3mm, left=4mm, right=4mm]
    \small
    \begin{itemize}[leftmargin=1.5em, itemsep=0.25em, label=---]
        \item Le sujet comporte \textbf{\pageref{LastPage} pages} numérotées de 1 à \pageref{LastPage}.
        \item Le candidat doit traiter les \textbf{4 exercices}.
        \item La qualité de la rédaction, la clarté et la précision des raisonnements seront prises en compte dans l'appréciation des copies.
        \item Tous les résultats doivent être justifiés, sauf mention contraire explicite.
        \item L'usage de tout document, téléphone ou appareil électronique connecté autre que la calculatrice autorisée est strictement interdit.
    \end{itemize}
\end{tcolorbox}

% -------------------------------------------------------------------------
% PAGE 2 : EXERCICE 1 (4 points)
% -------------------------------------------------------------------------
\newpage

\exercice[4]{1}{Thème : \ph{Thème de l'exercice 1}}

Pour chacune des affirmations suivantes, indiquer si elle est vraie ou fausse. Chaque réponse devra être soigneusement justifiée ; une réponse non justifiée ne rapporte aucun point. Une absence de réponse ou une réponse inexacte n'enlève pas de point. Les questions de cet exercice sont indépendantes.

\vspace{0.55cm}

\begin{enumerate}
    \item \ph{Contexte de l'affirmation 1}
    \[
        \ph{formule centrée}\,.
    \]
    \ph{Complément de contexte de l'affirmation 1}\\[0.2cm]
    \textbf{Affirmation 1 :} \ph{Énoncé de l'affirmation 1}
    
    \vspace{0.75cm}
    \item \ph{Contexte de l'affirmation 2}\\[0.2cm]
    \textbf{Affirmation 2 :} \ph{Énoncé de l'affirmation 2}
    
    \vspace{0.75cm}
    \item \ph{Contexte de l'affirmation 3}\\[0.2cm]
    \textbf{Affirmation 3 :} \ph{Énoncé de l'affirmation 3}
    
    \vspace{0.75cm}
    \item \ph{Contexte de l'affirmation 4}
    \[
        \ph{formule centrée}.
    \]
    \ph{Précision ou propriété admise}\\[0.2cm]
    \textbf{Affirmation 4 :} \ph{Énoncé de l'affirmation 4}
\end{enumerate}

% -------------------------------------------------------------------------
% PAGE 3 : EXERCICE 2 (5 points)
% -------------------------------------------------------------------------
\newpage

\exercice[5]{2}{Thème : \ph{Thème de l'exercice 2}}

\ph{Contexte de l'exercice 2}

\vspace{0.25cm}
\ph{Phrase d'introduction des notations} On note :
\begin{itemize}[label=\textbullet, itemsep=0.25em, topsep=0.15cm]
    \item $\ph{notation 1}$ \ph{définition 1} ;
    \item $\ph{notation 2}$ \ph{définition 2} ;
    \item $\ph{notation 3}$ \ph{définition 3}.
\end{itemize}

\vspace{0.45cm}
\noindent\textbf{\large Partie A --- \ph{Titre de la partie A}}
\vspace{0.25cm}

\begin{enumerate}
    \item \ph{Question A.1}
    
    \vspace{0.35cm}
    \item \ph{Question A.2}
    
    \vspace{0.35cm}
    \item \ph{Question A.3}
    
    \vspace{0.35cm}
    \item \ph{Question A.4} \textit{(\ph{consigne éventuelle}).}
\end{enumerate}

\vspace{0.55cm}
\noindent\textbf{\large Partie B --- \ph{Titre de la partie B}}
\vspace{0.25cm}

\begin{enumerate}
    \item \ph{Contexte de la question B.1}
    \begin{enumerate}[topsep=0.15cm, itemsep=0.3em]
        \item \ph{Question B.1.a}
        \item \ph{Question B.1.b}
        \item \ph{Question B.1.c} \textit{(\ph{consigne éventuelle}).}
    \end{enumerate}
    
    \vspace{0.45cm}
    \item \ph{Contexte de la question B.2}
    \begin{enumerate}[topsep=0.15cm, itemsep=0.3em]
        \item \ph{Question B.2.a}
        \item \ph{Question B.2.b}
    \end{enumerate}
\end{enumerate}

% -------------------------------------------------------------------------
% PAGE 4 : EXERCICE 3 (6 points)
% -------------------------------------------------------------------------
\newpage

\exercice[6]{3}{Thème : \ph{Thème de l'exercice 3}}

\ph{Contexte de l'exercice 3}
\[
    \ph{formule centrée}\,.
\]

\partie{A --- \ph{Titre de la partie A}}

\begin{enumerate}
    \item \ph{Question A.1}
    
    \item \ph{Question A.2}\\[0.05cm]
    \textit{(\ph{indication ou propriété admise})}.
    
    \item 
    \begin{enumerate}
        \item \ph{Question A.3.a}
        \[
            \ph{formule centrée}\,.
        \]
        \item \ph{Question A.3.b}
    \end{enumerate}
    
    \item \ph{Question A.4}
\end{enumerate}

\partie{B --- \ph{Titre de la partie B}}

\ph{Introduction de la partie B}
\[
    \ph{formule centrée}\,.
\]

\begin{enumerate}
    \item \ph{Question B.1}
    
    \item \ph{Question B.2}
    
    \item 
    \begin{enumerate}
        \item \ph{Question B.3.a}
        \item \ph{Question B.3.b}
        \item \ph{Question B.3.c}
    \end{enumerate}
    
    \item \ph{Question B.4}
    
    \item \ph{Question B.5 : introduction}
    
    \par\vspace{0.12cm}
    \noindent
    \begin{minipage}[t]{0.66\linewidth}
        \vspace{0pt}
        \begin{enumerate}[topsep=0pt, partopsep=0pt, leftmargin=1.5em, itemsep=0.45em]
            \item \ph{Question B.5.a}
            \item \ph{Question B.5.b}
        \end{enumerate}
    \end{minipage}\hfill
    \begin{minipage}[t]{0.32\linewidth}
        \vspace{0pt}
        \centering
        \begin{tcolorbox}[
            colback=gray!4,
            colframe=black!85,
            boxrule=0.6pt,
            arc=0.8mm,
            top=1.6mm,
            bottom=1.6mm,
            left=2.5mm,
            right=2mm,
            title={\footnotesize\textbf{\textsf{\ph{Titre de l'encadré}}}},
            coltitle=black,
            colbacktitle=gray!20
        ]
            \ttfamily\small
            \renewcommand{\arraystretch}{1.18}
            \begin{tabular}{r@{\hspace{0.45em}}|@{\hspace{0.45em}}l}
                1 & \ph{l.\,1} \\
                2 & \quad \ph{l.\,2} \\
                3 & \quad \ph{l.\,3} \\
                4 & \quad \ph{l.\,4} \makebox[1.4cm]{\dotfill} \\
                5 & \qquad \ph{l.\,5} \\
                6 & \qquad \ph{l.\,6} \makebox[1.4cm]{\dotfill} \\
                7 & \quad \ph{l.\,7}
            \end{tabular}
        \end{tcolorbox}
    \end{minipage}
\end{enumerate}

% -------------------------------------------------------------------------
% PAGE 5 : EXERCICE 4 (5 points)
% -------------------------------------------------------------------------
\newpage

\exercice[5]{4}{Thème : \ph{Thème de l'exercice 4}}

\ph{Introduction de l'exercice 4}
\[
    \ph{formule centrée}\,.
\]

\vspace{0.45cm}
\noindent\textbf{\large Partie A --- \ph{Titre de la partie A}}
\vspace{0.25cm}

\ph{Introduction de la partie A}
\[
    \ph{formule centrée}.
\]

\vspace{0.25cm}
\begin{enumerate}
    \item \ph{Question A.1}
    
    \vspace{0.35cm}
    \item \ph{Question A.2}
    
    \vspace{0.35cm}
    \item \ph{Question A.3}
    
    \vspace{0.35cm}
    \item \ph{Question A.4}
\end{enumerate}

\vspace{0.55cm}
\noindent\textbf{\large Partie B --- \ph{Titre de la partie B}}
\vspace{0.25cm}

\begin{enumerate}
    \item \ph{Question B.1}
    \[
        \ph{formule centrée}
    \]
    
    \vspace{0.30cm}
    \item \ph{Question B.2}
    
    \vspace{0.30cm}
    \item \ph{Question B.3}
    
    \vspace{0.30cm}
    \item \ph{Question B.4}
    
    \vspace{0.30cm}
    \item 
    \begin{enumerate}[topsep=0.15cm, itemsep=0.3em]
        \item \ph{Question B.5.a}
        \item \ph{Question B.5.b}
    \end{enumerate}
\end{enumerate}

\end{document}